\documentclass[11pt,oneside]{article}
\newcommand{\TopicTitle}{Apuntes del Tema 1: Foundations of Machine Learning — Classical Models \& Clustering}
\usepackage[a4paper,margin=2.6cm,headheight=15pt]{geometry}
\usepackage{fontspec}
\usepackage{polyglossia}
\setdefaultlanguage{spanish}
\setotherlanguage{english}
\setmainfont{Latin Modern Roman}
\setsansfont{Latin Modern Sans}
\setmonofont{Latin Modern Mono}[Scale=MatchLowercase]
\usepackage{amsmath,amssymb,mathtools}
\usepackage{graphicx}
\usepackage{xcolor}
\usepackage{booktabs}
\usepackage{tabularx}
\usepackage{array}
\usepackage{longtable}
\usepackage{enumitem}
\usepackage{listings}
\usepackage{fancyhdr}
\usepackage{titlesec}
\usepackage{microtype}
\usepackage{hyperref}
\usepackage{bookmark}
\usepackage{caption}

\definecolor{ink}{HTML}{243447}
\definecolor{accent}{HTML}{176B87}
\definecolor{codebg}{HTML}{F3F5F7}
\definecolor{codegray}{HTML}{65727E}

\hypersetup{
  colorlinks=true,
  linkcolor=accent,
  urlcolor=accent,
  citecolor=accent,
  pdftitle={\TopicTitle},
  pdfauthor={Foundations of Machine Learning and Deep Learning}
}
\setcounter{secnumdepth}{3}
\setcounter{tocdepth}{3}
\setlist[itemize]{leftmargin=1.7em,itemsep=2pt,topsep=4pt}
\setlist[enumerate]{leftmargin=1.9em,itemsep=2pt,topsep=4pt}
\setlength{\parindent}{0pt}
\setlength{\parskip}{5pt plus 1pt minus 1pt}
\titleformat{\section}{\Large\sffamily\bfseries\color{ink}}{\thesection}{0.8em}{}
\titleformat{\subsection}{\large\sffamily\bfseries\color{accent}}{\thesubsection}{0.8em}{}
\titleformat{\subsubsection}{\normalsize\sffamily\bfseries\color{ink}}{\thesubsubsection}{0.8em}{}
\titlespacing*{\section}{0pt}{2.0em}{0.7em}
\titlespacing*{\subsection}{0pt}{1.4em}{0.4em}
\titlespacing*{\subsubsection}{0pt}{1.0em}{0.3em}
\lstset{
  backgroundcolor=\color{codebg},
  basicstyle=\ttfamily\small,
  breaklines=true,
  breakatwhitespace=false,
  columns=fullflexible,
  frame=single,
  framerule=0.3pt,
  rulecolor=\color{codegray},
  numbers=left,
  numbersep=8pt,
  numberstyle=\scriptsize\color{codegray},
  showstringspaces=false,
  keepspaces=true,
  tabsize=4,
  keywordstyle=\color{accent}\bfseries,
  commentstyle=\color{codegray}\itshape,
  stringstyle=\color{purple}
}
\captionsetup{font=small,labelfont=bf}
\pagestyle{fancy}
\fancyhf{}
\fancyhead[L]{\small\sffamily\color{codegray}Foundations of Machine Learning and Deep Learning}
\fancyhead[R]{\small\sffamily\color{codegray}\nouppercase{\rightmark}}
\fancyfoot[C]{\thepage}
\renewcommand{\headrulewidth}{0.4pt}
\renewcommand{\footrulewidth}{0pt}
\newcommand{\topiccover}[1]{%
  \begin{titlepage}
    \centering
    \vspace*{2.5cm}
    {\Large\sffamily\color{accent}APUNTES UNIVERSITARIOS\par}
    \vspace{1.5cm}
    {\Huge\sffamily\bfseries\color{ink}Foundations of Machine Learning\\[0.25em]
    and Deep Learning\par}
    \vspace{1.8cm}
    {\rule{0.72\textwidth}{1pt}\par}
    \vspace{1.2cm}
    {\LARGE\sffamily\bfseries #1\par}
    \vfill
    {\large\sffamily\color{codegray}\today\par}
  \end{titlepage}
}

\begin{document}
\topiccover{Apuntes del Tema 1: Foundations of Machine Learning — Classical Models \& Clustering}
\tableofcontents
\clearpage
\markboth{Apuntes del Tema 1: Foundations of Machine Learning — Classical Models \& Clustering}{Apuntes del Tema 1: Foundations of Machine Learning — Classical Models \& Clustering}
\section{Apuntes del Tema 1: Foundations of Machine Learning — Classical Models \& Clustering}

\subsection{1. Conceptos Fundamentales y Paradigmas de Aprendizaje}

\subsubsection{1.1 Trazabilidad y Definición}

\begin{itemize}

\item \textbf{Machine Learning:} Disciplina enfocada en la creación de sistemas capaces de realizar tareas mediante el reconocimiento de patrones y la toma de decisiones basadas en datos de experiencia (ejemplos), evaluando el desempeño mediante la capacidad de \textbf{generalización} hacia nuevos datos no vistos y previniendo el \emph{overfitting}.

\item \textbf{Estructura de Datos Básica:}

\item \textbf{Dataset:} Conjunto de observaciones \(\mathcal{D} = \{(x^{(i)}, y^{(i)})\}_{i=1}^{n}\).

\item \textbf{Instancia (\(x\)):} Vector de características de dimensión \(d\), donde \(x \in \mathbb{R}^d\).

\item \textbf{Etiqueta/Target (\(y\)):} Valor objetivo o clase a predecir.

\item \textbf{Modelo Paramétrico:} Función de mapeo \(f(x; \theta) \rightarrow \hat{y}\) parametrizada por pesos y sesgos \(\theta = (w, b)\).

\end{itemize}

\subsubsection{1.2 Paradigmas de Aprendizaje}

\begin{itemize}

\item \textbf{Supervisado:} Aprendizaje mediante pares entrada-salida etiquetados \((x, y)\). Se divide en:

\item \textbf{Regresión:} El target \(y\) es un valor continuo.

\item \textbf{Clasificación:} El target \(y\) representa una categoría discreta.

\item \textbf{No Supervisado:} Descubrimiento de estructuras implícitas o agrupamientos en datos sin etiquetas \(x\) (ej. \emph{Clustering}, reducción de dimensionalidad).

\item \textbf{Otros Paradigmas Relevantes:} \emph{Semi-supervised}, \emph{Self-supervised} (pretext tasks), \emph{Transfer Learning} (reutilización de backbones preentrenados), \emph{Active Learning}, \emph{Online/Incremental Learning} y \emph{Federated Learning}.

\end{itemize}

\subsection{2. k-Nearest Neighbors (kNN)}

\subsubsection{2.1 Principio e Intuición}

\begin{itemize}

\item \textbf{Algoritmo Lazy Learner:} No requiere una fase de entrenamiento o ajuste de parámetros explícito (\(f(x)\) se evalúa en tiempo de consulta).

\item \textbf{Regla de Clasificación:} Asigna la etiqueta a una muestra de prueba \(x_{\text{query}}\) según el voto mayoritario (o promedio ponderado por la distancia inversa) de sus \(k\) vecinos más cercanos en el espacio métrico.

\end{itemize}

\subsubsection{2.2 Métricas de Distancia}

\begin{itemize}

\item \textbf{Distancia Euclídea (\(L_2\)):}

\end{itemize}

\[\text{dist}_{\text{Euclidean}}(x, x') = \sqrt{\sum_{j=1}^{d} (x_j - x_j')^2}\]

\begin{itemize}

\item \textbf{Distancia Manhattan (\(L_1\)):}

\end{itemize}

\[\text{dist}_{\text{Manhattan}}(x, x') = \sum_{j=1}^{d} \vert{}x_j - x_j'\vert{}\]

\begin{itemize}

\item \textbf{Minkowski (\(L_p\)):} 

\end{itemize}

\[\text{dist}_{L^p}(x, x') = \left( \sum_{j=1}^{d} \vert{}x_j - x_j'\vert{}^p \right)^{1/p}\]

\subsubsection{2.3 Selección de \(k\) y Efectos de Escala}

\begin{itemize}

\item \textbf{Efecto de \(k\):}

\item \(k=1\) (Bajo): Frontera extremadamente flexible, muy sensible al ruido y propenso a sobreajuste (\emph{overfitting}).

\item \(k\) elevado: Fronteras muy suaves, riesgo de subajuste (\emph{underfitting}) al mezclar vecindarios.

\item \textbf{Necesidad de Normalización:} Dado que kNN depende de distancias geométricas, las variables con rangos numéricos elevados dominan el cálculo. Es indispensable estandarizar las características (\(x' = \frac{x - \mu}{\sigma}\)) antes de calcular las distancias.

\item **La Maldición de la Dimensionalidad (\emph{Curse of Dimensionality}):** A medida que la dimensión \(d\) del espacio crece, el volumen aumenta exponencialmente, volviendo los datos muy dispersos y reduciendo la capacidad discriminativa de las distancias geométricas entre los puntos.

\end{itemize}

\subsection{3. Naive Bayes}

\subsubsection{3.1 Fundamento Probabilístico Generativo}

\begin{itemize}

\item \textbf{Modelo Generativo:} Modela la probabilidad de la clase a priori \(P(y)\) y la verosimilitud de los datos \(P(x\vert{}y)\) para estimar la probabilidad a posteriori \(P(y\vert{}x)\) aplicando el Teorema de Bayes:

\end{itemize}

\[P(y\vert{}x) = \frac{P(x\vert{}y)P(y)}{P(x)} \propto P(y) \prod_{j=1}^{d} P(x_j\vert{}y)\]

\begin{itemize}

\item \textbf{Regla MAP (Maximum A Posteriori):} Selección de la clase que maximiza el producto: 

\end{itemize}

\[\hat{y} = \arg\max_{y} P(y) \prod_{j=1}^{d} P(x_j\vert{}y)\]

\begin{itemize}

\item \textbf{Asunción "Naive" (Ingenua):} Presupone la \textbf{independencia condicional} de todas las características \(x_j\) dada la clase \(y\).

\end{itemize}

\subsubsection{3.2 Variantes Principales de Naive Bayes}

\begin{center}
\begin{tabularx}{\textwidth}{>{\raggedright\arraybackslash}X >{\raggedright\arraybackslash}X >{\raggedright\arraybackslash}X}
\toprule
Variante & Tipo de Variable (\(x_j\)) & Estimación de Verosimilitud \(P(x_j\vert{}y)\) \\
\midrule
\textbf{Categorical NB} & Categóricas/Nominales discretas & \(P(x_j = v \mid y) = \frac{\text{count}(x_j=v, y) + \alpha}{\text{count}(y) + \alpha \cdot K_j}\) \\
\textbf{Multinomial NB} & Conteos discretos (frecuencias de texto/Bag-of-Words) & \(P(x_j \mid y) = \frac{\text{count}(x_j \text{ en clase } y) + \alpha}{\sum_{k} \text{count}(x_k \text{ en clase } y) + \alpha \cdot V}\) \\
\textbf{Gaussian NB} & Características continuas & \(P(x_j \mid y) = \frac{1}{\sqrt{2\pi \sigma_{j,y}^2}} \exp\left(-\frac{(x_j - \mu_{j,y})^2}{2\sigma_{j,y}^2}\right)\) \\
\bottomrule
\end{tabularx}
\end{center}

\begin{itemize}

\item \textbf{Laplace Smoothing (\(\alpha=1\)):} Corrige el problema de la \textbf{probabilidad cero} (frecuencia nula en datos de entrenamiento que anularía todo el producto de verosimilitudes) sumando el parámetro \(\alpha\) en el numerador y \(\alpha \cdot K_j\) (o \(\alpha \cdot V\)) en el denominador.

\end{itemize}

\subsection{4. Árboles de Decisión (Decision Trees)}

\subsubsection{4.1 Concepto y Particionado Recursivo}

\begin{itemize}

\item Modelos no paramétricos e interpretables para clasificación y regresión que dividen recursivamente el espacio de características mediante reglas ortogonales a los ejes.

\end{itemize}

\subsubsection{4.2 Comparativa de Algoritmos y Criterios de División}

\begin{enumerate}

\item \textbf{ID3 (Iterative Dichotomiser 3):}

\end{enumerate}

\begin{itemize}

\item Solo admite atributos categóricos y no realiza poda.

\item Utiliza la \textbf{Entropía} \(H(S)\) y la \textbf{Ganancia de Información} \(IG(S, A)\):

\end{itemize}

\texttt{\$\$H(S) = - \textbackslash\{\}sum\_\{i=1\}\textasciicircum{}\{c\} p\_i \textbackslash\{\}log\_2(p\_i)\$\$}\\

\texttt{\$\$IG(S, A) = H(S) - \textbackslash\{\}sum\_\{v \textbackslash\{\}in \textbackslash\{\}text\{Valores\}(A)\} \textbackslash\{\}frac\{\textbackslash\{\}vert\{\}S\_v\textbackslash\{\}vert\{\}\}\{\textbackslash\{\}vert\{\}S\textbackslash\{\}vert\{\}\} H(S\_v)\$\$}\\

\begin{enumerate}

\item \textbf{C4.5:}

\end{enumerate}

\begin{itemize}

\item Soporta variables continuas mediante la evaluación de umbrales binarios de corte \(t\).

\item Sustituye el \(IG\) por el \textbf{Gain Ratio} para corregir el sesgo de ID3 hacia atributos con excesivos valores únicos (como un \texttt{ID\_Cliente}):

\end{itemize}

\texttt{\$\$\textbackslash\{\}text\{GainRatio\}(S, A) = \textbackslash\{\}frac\{IG(S, A)\}\{\textbackslash\{\}text\{SplitInfo\}(A)\}\$\$}\\

\texttt{\$\$\textbackslash\{\}text\{SplitInfo\}(A) = -\textbackslash\{\}sum\_\{i=1\}\textasciicircum{}\{k\} \textbackslash\{\}frac\{\textbackslash\{\}vert\{\}S\_i\textbackslash\{\}vert\{\}\}\{\textbackslash\{\}vert\{\}S\textbackslash\{\}vert\{\}\} \textbackslash\{\}log\_2\textbackslash\{\}left(\textbackslash\{\}frac\{\textbackslash\{\}vert\{\}S\_i\textbackslash\{\}vert\{\}\}\{\textbackslash\{\}vert\{\}S\textbackslash\{\}vert\{\}\}\textbackslash\{\}right)\$\$}\\

\begin{itemize}

\item Aplica poda pesimista posterior al entrenamiento.

\end{itemize}

\begin{enumerate}

\item \textbf{CART (Classification and Regression Trees):}

\end{enumerate}

\begin{itemize}

\item Realiza exclusivamente divisiones binarias.

\item Minimiza la \textbf{Impureza de Gini} en clasificación:

\end{itemize}

\texttt{\$\$G(S) = 1 - \textbackslash\{\}sum\_\{i=1\}\textasciicircum{}\{c\} p\_i\textasciicircum{}2\$\$}\\

\begin{itemize}

\item En regresión, minimiza la varianza/suma de errores cuadráticos (SSE).

\end{itemize}

\subsection{5. Clustering (Agrupamiento No Supervisado)}

\subsubsection{5.1 k-means}

\begin{itemize}

\item \textbf{Objetivo:} Particionar \(X\) en \(c\) clusters disjuntos \(\omega_1, \dots, \omega_c\) minimizando la varianza dentro de cada cluster (suma de distancias al cuadrado al centroide \(\mu_i\)).

\item \textbf{Paso a Paso del Algoritmo:}

\end{itemize}

\begin{enumerate}

\item Inicialización aleatoria de \(k\) centroides.

\item Asignación de cada punto al centroide más cercano.

\item Reordenamiento y recalculo del punto medio/centroide de cada cluster.

\item Repetición iterativa hasta convergencia.

\end{enumerate}

\begin{itemize}

\item \textbf{Selección de \(k\):} Se evalúa empíricamente mediante el \textbf{Elbow Method} (buscando el punto de inflexión en la curva de la inercia/varianza intra-cluster).

\end{itemize}

\subsubsection{5.2 Clustering Jerárquico Aglomerativo (Agglomerative)}

\begin{itemize}

\item \textbf{Mecanismo:} Construye una estructura en árbol (\emph{Dendrograma}) fusionando sucesivamente los dos clusters más cercanos desde un estado inicial donde cada punto es un cluster individual.

\item \textbf{Criterios de Enlace (Linkage Criteria):}

\item \textbf{Single Linkage:} Distancia mínima entre miembros de dos clusters (\(d_{\min}(C_1, C_2)\)).

\item \textbf{Complete Linkage:} Distancia máxima entre miembros de dos clusters (\(d_{\max}(C_1, C_2)\)).

\item \textbf{Average Linkage:} Distancia promedio entre todos los pares de elementos entre clusters.

\end{itemize}

\subsection{6. Práctica y Aplicación en Código PyTorch / Scikit-Learn}

\begin{lstlisting}[language=Python]

import numpy as np

import torch

import torch.nn as nn

import torch.optim as optim

from sklearn.neighbors import KNeighborsClassifier

from sklearn.naive_bayes import GaussianNB

from sklearn.tree import DecisionTreeClassifier

from sklearn.cluster import KMeans

from sklearn.preprocessing import StandardScaler



# --------------------------------------------------------------------------

# 1. Modelos Clásicos en Scikit-Learn (kNN, Naive Bayes, Decision Trees)

# --------------------------------------------------------------------------

def train_classical_models(X_train, y_train, X_test):

    # kNN: requiere estandarización previa de características

    scaler = StandardScaler()

    X_train_scaled = scaler.fit_transform(X_train)

    X_test_scaled = scaler.transform(X_test)

    

    knn = KNeighborsClassifier(n_neighbors=5, metric='euclidean')

    knn.fit(X_train_scaled, y_train)

    

    # Gaussian Naive Bayes

    gnb = GaussianNB()

    gnb.fit(X_train, y_train)

    

    # Árbol de Decisión CART (Criterio Gini)

    tree = DecisionTreeClassifier(criterion='gini', max_depth=5)

    tree.fit(X_train, y_train)

    

    return knn, gnb, tree



# --------------------------------------------------------------------------

# 2. Clustering (k-means)

# --------------------------------------------------------------------------

def run_kmeans_clustering(X_data, n_clusters=3):

    kmeans = KMeans(n_clusters=n_clusters, random_state=42, n_init=10)

    cluster_labels = kmeans.fit_predict(X_data)

    centroids = kmeans.cluster_centers_

    return cluster_labels, centroids



# --------------------------------------------------------------------------

# 3. MLP Baseline con Vectorización (Aplanado) en PyTorch

# --------------------------------------------------------------------------

class MLPClassifier(nn.Module):

    """

    MLP Baseline para imágenes o datos vectorizados.

    Demuestra la pérdida de estructura espacial al aplanar un tensor (C, H, W) -> (C * H * W)

    """

    def __init__(self, input_shape=(3, 64, 64), num_classes=10, hidden_units=[512, 256], dropout_rate=0.3):

        super(MLPClassifier, self).__init__()

        

        # Dimensión de la entrada aplanada

        self.input_dim = input_shape[0] * input_shape[1] * input_shape[2]

        

        layers = []

        in_dim = self.input_dim

        

        for h_dim in hidden_units:

            layers.append(nn.Linear(in_dim, h_dim))

            layers.append(nn.ReLU())

            layers.append(nn.Dropout(p=dropout_rate))

            in_dim = h_dim

            

        # Capa lineal de salida (retorna logits para CrossEntropyLoss)

        layers.append(nn.Linear(in_dim, num_classes))

        

        self.network = nn.Sequential(*layers)



    def forward(self, x):

        # Aplanado explícito: (Batch, Channels, Height, Width) -> (Batch, C * H * W)

        x = x.view(x.size(0), -1)

        return self.network(x)



if __name__ == "__main__":

    device = torch.device("cuda" if torch.cuda.is_available() else "cpu")

    dummy_images = torch.randn(16, 3, 64, 64).to(device)

    model = MLPClassifier(input_shape=(3, 64, 64), num_classes=10).to(device)

    logits = model(dummy_images)

    print(f"Forma de la entrada : {dummy_images.shape}")

    print(f"Forma de los logits : {logits.shape}")

\end{lstlisting}
\end{document}
